\documentclass[a4paper,12pt]{article}

\usepackage[slovene]{babel}
\usepackage[T1]{fontenc}
\usepackage[utf8]{inputenc}
\usepackage{amsmath,amssymb}
\usepackage[margin=2.2cm,headheight=16pt]{geometry}
\usepackage{enumitem}
\usepackage{xcolor}
\usepackage{fancyhdr}
\usepackage{titlesec}

\definecolor{bordotemno}{RGB}{101,8,34}
\definecolor{bordosvetlo}{RGB}{153,67,88}
\pagestyle{fancy}
\fancyhf{}
\fancyhead[L]{\small Matematika \textbar\ Funkcije}
\fancyhead[R]{\small Rešitve preverjanja znanja}
\fancyfoot[C]{\thepage}
\renewcommand{\headrulewidth}{0.4pt}
\titleformat{\section}{\large\bfseries\color{bordotemno}}{}{0pt}{}
\titlespacing*{\section}{0pt}{1.3em}{0.7em}
\setlist[enumerate]{itemsep=5pt,topsep=4pt}

\begin{document}

\begin{center}
	{\LARGE\bfseries\color{bordotemno} REŠITVE PREVERJANJA ZNANJA}\par
	\vspace{2mm}
	{\large Funkcije \textbar\ Matematika, 2. letnik}
\end{center}

\noindent\textit{Rešitve vsebujejo možne postopke. Pri računskih nalogah je treba priznati tudi drugačen pravilen postopek.}

\section*{Naloga 1: Pojem funkcije \hfill\normalfont\normalsize(10 točk)}
\begin{enumerate}[label=\arabic*)]
	\item Prirejanje je funkcija, če vsakemu elementu $x$ iz definicijskega območja priredi natanko eno vrednost $y$. Dvema različnima argumentoma je lahko prirejena ista vrednost.
	\item Definicijsko območje $D_f$ je množica vseh dovoljenih argumentov $x$. Zaloga vrednosti $Z_f$ je množica vseh vrednosti $f(x)$, ki jih funkcija dejansko doseže.
	\item Funkcijo lahko podamo s predpisom, tabelo vrednosti, grafom, besednim opisom ali diagramom prirejanja. Zadoščajo poljubni trije načini.
\end{enumerate}
\textbf{Točkovanje:} 3 + 4 + 3 točke.

\newpage
\section*{Naloga 2: Lastnosti funkcij \hfill\normalfont\normalsize(10 točk)}
\begin{enumerate}[label=\arabic*)]
	\item Pred preverjanjem mora biti definicijsko območje simetrično glede na $0$. Nato velja:
	\[
		f(-x)=f(x)\quad\Longrightarrow\quad f\text{ je soda},\qquad
		f(-x)=-f(x)\quad\Longrightarrow\quad f\text{ je liha}.
	\]
	\item Graf sode funkcije je simetričen glede na os $y$, graf lihe funkcije pa glede na koordinatno izhodišče.
	\item Funkcija je injektivna, če različna argumenta vedno določita različni funkcijski vrednosti. Grafični preizkus: vsaka vodoravna premica seka graf največ enkrat.
	\item Funkcija ima inverzno funkcijo natanko tedaj, ko je injektivna na svojem definicijskem območju. Inverzna funkcija je definirana na zalogi vrednosti prvotne funkcije.
\end{enumerate}
\textbf{Točkovanje:} 4 + 2 + 2 + 2 točki.

\newpage
\section*{Naloga 3: Vrednosti, ničle in definicijsko območje \hfill\normalfont\normalsize(15 točk)}
\begin{enumerate}[label=\arabic*)]
	\item Za $f(x)=2x-3$:
	\[
		f(-1)=2(-1)-3=-5,\qquad f(0)=2\cdot0-3=-3.
	\]
	Ničlo dobimo iz $2x-3=0$, zato je $x=\dfrac32$.
	\item Pri $g(x)=\sqrt{5-2x}$ mora biti podkoren izraz nenegativen:
	\[
		5-2x\geq0\quad\Longrightarrow\quad x\leq\frac52,
		\qquad D_g=(-\infty,\tfrac52].
	\]
	Ničla je tam, kjer je $5-2x=0$, torej $x=\dfrac52$.
	\item Imenovalec ne sme biti nič:
	\[
		x^2-9\ne0\quad\Longrightarrow\quad (x-3)(x+3)\ne0.
	\]
	Zato je $D_r=\mathbb{R}\setminus\{-3,3\}$.
\end{enumerate}
\textbf{Točkovanje:} 6 + 5 + 4 točke.

\newpage
\section*{Naloga 4: Raziskovanje kvadratne funkcije \hfill\normalfont\normalsize(20 točk)}
Dana je $p(x)=x^2-4x+3$.
\begin{enumerate}[label=\arabic*)]
	\item Ker je polinom, je $D_p=\mathbb{R}$. Ničle:
	\[
		x^2-4x+3=(x-1)(x-3)=0\quad\Longrightarrow\quad x_1=1,\;x_2=3.
	\]
	\item Parabola je obrnjena navzgor. Zato je $p(x)>0$ za $x\in(-\infty,1)\cup(3,\infty)$, $p(x)<0$ za $x\in(1,3)$, pri $x=1$ in $x=3$ pa je $p(x)=0$.
	\item Dopolnimo do kvadrata:
	\[
		p(x)=x^2-4x+4-1=(x-2)^2-1.
	\]
	Temenska oblika je $(x-2)^2-1$, teme je $T(2,-1)$, najmanjša vrednost pa $-1$.
	\item Ker je teme pri $x=2$ in parabola odprta navzgor, funkcija pada na $(-\infty,2]$ ter narašča na $[2,\infty)$. Zaloga vrednosti je $Z_p=[-1,\infty)$.
\end{enumerate}
\textbf{Točkovanje:} 4 + 4 + 6 + 6 točk.

\newpage
\section*{Naloga 5: Premik in razteg grafa \hfill\normalfont\normalsize(20 točk)}
Dana je $h(x)=-2(x-1)^3+4$.
\begin{enumerate}[label=\arabic*)]
	\item Iz grafa $y=x^3$ najprej dobimo $y=(x-1)^3$ s premikom za $1$ v desno. Množenje z $-2$ pomeni navpični razteg s faktorjem $2$ in zrcaljenje čez os $x$. Prištevanje $4$ graf premakne za $4$ navzgor.
	\item Kubična funkcija ima definicijsko območje in zalogo vrednosti $\mathbb{R}$. Ti lastnosti ohranijo tudi premiki, razteg in zrcaljenje, zato $D_h=Z_h=\mathbb{R}$.
	\item Začetna vrednost:
	\[
		h(0)=-2(0-1)^3+4=6.
	\]
	Ničla:
	\[
		-2(x-1)^3+4=0\quad\Longrightarrow\quad (x-1)^3=2
		\quad\Longrightarrow\quad x=1+\sqrt[3]{2}.
	\]
	Presečišči z osema sta $(0,6)$ in $(1+\sqrt[3]{2},0)$.
	\item Za skico označimo še prevoj $T(1,4)$ in uporabimo značilne točke, na primer:
	\[
		h(-1)=20,\quad h(0)=6,\quad h(1)=4,\quad h(2)=2,\quad h(3)=-12.
	\]
	Graf je padajoča kubična krivulja v obliki črke S; poteka skozi navedene točke, zlasti skozi presečišči z osema in prevoj.
\end{enumerate}
\textbf{Točkovanje:} 5 + 2 + 5 + 8 točk.

\newpage
\section*{Naloga 6: Koreni in racionalne funkcije \hfill\normalfont\normalsize(25 točk)}
\begin{enumerate}[label=\arabic*)]
	\item Pogoj za obstoj korena je $2x+3\geq0$. Ker je desna stran enačbe koren, mora veljati tudi $x\geq0$. Enačbo kvadriramo:
	\[
		\sqrt{2x+3}=x\quad\Longrightarrow\quad 2x+3=x^2
		\quad\Longrightarrow\quad (x-3)(x+1)=0.
	\]
	Kandidata sta $x=3$ in $x=-1$, vendar pogoj $x\geq0$ izpolnjuje le $x=3$. Preverjanje: $\sqrt{2\cdot3+3}=\sqrt9=3$. Rešitev je $\boxed{x=3}$.
	\item Dana je $q(x)=\dfrac{1}{x^2-4}$.
	\begin{enumerate}[label=\alph*)]
		\item Imenovalec ne sme biti nič, zato
		\[
			D_q=\mathbb{R}\setminus\{-2,2\}.
		\]
		Navpični asimptoti sta $x=-2$ in $x=2$, vodoravna asimptota pa $y=0$.
		Za zalogo vrednosti naj bo $u=x^2\geq0$, pri čemer $u\ne4$. Če je $0\leq u<4$, potem $q=1/(u-4)\in(-\infty,-1/4]$; če je $u>4$, potem $q\in(0,\infty)$. Torej
		\[
			Z_q=(-\infty,-\tfrac14]\cup(0,\infty).
		\]
		\item Območje je simetrično glede na $0$ in
		\[
			q(-x)=\frac{1}{(-x)^2-4}=\frac{1}{x^2-4}=q(x),
		\]
		zato je $q$ soda.
		\item Rešimo neenačbo:
		\[
			\frac{1}{x^2-4}\geq\frac15
			\quad\Longleftrightarrow\quad
			\frac{9-x^2}{5(x^2-4)}\geq0.
		\]
		Kritične vrednosti so $-3,-2,2,3$. Vrednosti $-2$ in $2$ izločimo iz definicijskega območja; pri $-3$ in $3$ je izraz enak nič. Predznak izraza je nenegativen na $[-3,-2)$ in $(2,3]$. Zato je rešitev
		\[
			\boxed{x\in[-3,-2)\cup(2,3]}.
		\]
	\end{enumerate}
\end{enumerate}
\textbf{Točkovanje:} 10 + 8 + 2 + 5 točk.

\end{document}
